\section{Introduction}
\label{sec:intro}

Federated LoRA methods for instruction tuning report method differences that are often small relative to the noise in a single evaluation. FFA-LoRA~\cite{sun2024ffalora} publishes means over $20$ runs with standard deviations; in the non-private setting the LoRA standard deviation on MNLI-matched ($10.7$) exceeds the mean gap to FFA-LoRA ($3.02$ points). FLoRA~\cite{wang2024flora} compares to FedIT~\cite{zhang2024fedit} with single reported MMLU and MT-bench values after $1$ to $3$ communication rounds, without reported seed replication, including a Llama/Alpaca MMLU margin of $0.44$ points that the paper itself calls ``marginal.'' Those practices establish the motivating fact: published federated LoRA contrasts can sit inside, or only slightly above, ordinary run-to-run variation. What they do not establish is how much of that variation comes from the non-IID partition draw itself.

In Dirichlet or category-based federated simulations, the ``non-IID'' condition is commonly instantiated from a single partition seed: the Dirichlet construction itself~\cite{hsu2019measuring} and the partition catalogues built on it~\cite{li2022niidbench} treat the draw as a setting rather than a source of variance. In the evaluations cited above, training seeds are sometimes repeated while the partition seed is not. If held-out loss moves with the partition draw by as much as it moves with the training seed, then a ranking obtained from one or two partitions is not a property of the methods alone. The open question is quantitative: how large is the partition component, how often do few-draw rankings flip, and how many draws are needed for a target effect size?

We answer three pre-registered questions on FedIT, FFA-LoRA, and FLoRA, using held-out next-token loss on Dolly as the response. All differences are in held-out next-token loss, where the unadapted base models score $2.121$ and $2.169$ and fine-tuned runs fall between $1.69$ and $1.841$. The design freezes the federation (ten clients, full participation, fifteen rounds, one local epoch), the LoRA configuration, and the analysis plan before any production run, and separates partition seeds from training seeds at two model scales.

\begin{itemize}
\item \textbf{RQ1.} At Dirichlet $\alpha{=}0.1$, what share of the variance comes from the partition draw, the partition-by-method interaction, and the training seed?
\item \textbf{RQ2.} How often does an evaluation using $k$ partition draws rank the three methods differently from the full-design ranking?
\item \textbf{RQ3.} How many partition draws are needed to detect held-out loss differences of $0.005$, $0.01$, $0.02$, and $0.05$ at $80\%$ power, for paired and unpaired designs?
\end{itemize}

An exploratory RQ4 on partition statistics is reported but not used for claims.

\paragraph{Advance over the state of the art.}
Prior federated LoRA evaluations either report multi-seed means without separating the partition draw from the training seed, or treat the non-IID split as a fixed setting rather than a variance component. Adjacent work decomposes variance over training or unlearning seeds and over benchmark variants, but not over the federated partition draw for LoRA instruction tuning. This paper's advance is to isolate that draw in a pre-registered design, quantify its share relative to interaction and residual noise, and turn the measurement into an explicit reporting protocol that states how many draws were used, that comparisons were paired on those draws, and what minimum difference the design can detect. The same design also shows the cost of too few draws: extending LLaMA-3.2-3B from six to ten partitions reverses the FedIT--FLoRA top-two order and grows the interaction component by a factor of $2.66$.

\paragraph{Contributions.}
(1)~A pre-registered $120$-run design at two scales (TinyLlama-1.1B-Chat and LLaMA-3.2-3B), extended by $84$ further runs, that separates partition seeds from training seeds under a frozen federation and LoRA configuration. (2)~A method-of-moments variance decomposition with cluster-bootstrap intervals, showing on TinyLlama that partition and interaction shares dominate the residual ($0.508$ and $0.469$ versus $0.023$), while on LLaMA-3.2-3B at $p{=}10$ the interaction dominates ($0.992$) and the partition main-effect interval includes zero after truncation. (3)~A ranking-stability analysis conditioned on effect size: the near-tied FedIT--FLoRA pair flips under few draws (three-draw paired flip probability $0.377$ on TinyLlama and $0.493$ on LLaMA-3.2-3B), whereas gaps of about $0.023$ to $0.031$ never reverse in a single draw. (4)~A draws-needed table for pre-registered deltas and for the observed gaps, including $206$ and more than $1000$ paired draws for the near-tie at the two scales, and the observation that where the interaction dominates, pairing on partitions buys little. (5)~The Partition-Draw Reporting Protocol as three numbered steps (partition-draw count and seeds; paired comparisons; detectable minimum difference). We make no method recommendation.

Section~\ref{sec:related} situates the work against benchmark-variance studies, non-IID simulation practice, and federated LoRA evaluation. Section~\ref{sec:setup} states the design and variance model. Section~\ref{sec:results} reports the measurements. Section~\ref{sec:discussion} gives reporting recommendations and limitations. Section~\ref{sec:conclusion} concludes.
